\section{Extending \MMCC{} to ``$b$-min-sep Sampling''}\label{sec:mmcc-extension}

\cite{BandedMF} analyzed the $b$-banded matrix mechanism under the following scheme, which we'll call ``$b$-min-sep sampling'': We partition the dataset $D$ into $b$ equal-size subsets, $D_1, D_2, \ldots D_b$. To compute $\bfx_i$, we use independently include each element of $D_{i \pmod b}$ (where we say $i \pmod b = b$ if $b$ divides $i$) with probability $bp$; here, we write the sampling probability in these rounds as $bp$ instead of $p$ to reflect the fact that the average example still participates in fraction $p$ of rounds in expectation for any choice of $b$.

We give a generalization of \MMCC{} that analyzes the matrix mechanism under $b$-min-sep sampling, that matches the analysis of \cite{BandedMF} when $\bfC$ is $b$-banded but can generalize to arbitrary lower triangular matrices. In other words, this generalization of \MMCC{} subsumes the analysis in \cite{BandedMF}. 


\begin{figure}
\begin{algorithm}[H]
\caption{\texttt{Generalized-}\MMCC{}}
\begin{algorithmic}[1]
\State \textbf{Input:} Matrix $\bfC$, sampling probability $p$, noise standard deviation $\sigma$, probabilities $\delta_1, \delta_2$, min-sep $b$.
\State Delete all columns of $\bfC$ except columns $1, b + 1, 2b + 1 \ldots$
\State $\{\tilde{p}_{i,j}\}_{i \in [n], j \in [\lceil n / b \rceil} \leftarrow$\texttt{GeneralizedProbabilityTailBounds}($\bfC, bp, \sigma, b  \delta_1)$.\newline\Comment{$\tilde{p}_{i,j}$ is a high-probability upper bound on the probability that an example participated in round $j$, conditioned on output in rounds $1$ to $i-1$.}
\State $\tilde{p}^{(b)}_{i,j} = \tilde{p}_{(i-1)b + 1, (j-1)b + 1}$
\State $\bfC^{(b)}_{i,j} = \ltwo{\bfC_{(i-1)b + 1 : ib, (j-1)b + 1}}$
\For{$i \in [\lceil n/b \rceil]$} 
\State $PLD_i \leftarrow$ PLD of $\calM_{PMoG}(\{\tilde{p}^{(b)}_{i, j}\}_{j \in [\lceil n/b \rceil]}, \{\bfC^{(b)}_{i, j}\}_{j \in [\lceil n/b \rceil]}).$
\EndFor
\State $PLD \leftarrow$ convolution of $\{PLD_i\}_{i \in [\lceil n/b \rceil]}$.
\State \Return $\min\left(\{\epsilon: PLD\text{ satisfies }(\epsilon, \delta_2)\text{-DP}\}\right)$.
\end{algorithmic}
\end{algorithm}
\caption{Extension of \CPTB{} to $b$-min-sep sampling.}
\label{fig:cptb-extension}
\end{figure}

\begin{figure}
\begin{algorithm}[H]
\caption{\texttt{GeneralizedProbabilityTailBounds}($\bfC, p, \sigma, \delta_1)$}
\begin{algorithmic}[1]
\State \textbf{Input:} Matrix $\bfC \in \mathbb{R}^{m \times n}$, sampling probability $p$, noise standard deviation $\sigma$, probability $\delta_1$.
\State $\delta'$ = $\frac{\delta_1}{2 \cdot (nnz(\bfC) - n)}$ \Comment{$nnz$ is the number of non-zeros.}
\State $z = \Phi^{-1}(1 - \delta')$ \Comment{Tail bound on normal distribution; here, $\Phi$ is the standard normal CDF.}
\For{$i \in [m], j \in [n]$}
\If{$\bfC_{i, j} = 0$}
\State $\tilde{p}_{i, j} = 1$
\Else 
\State $s_{i,j} = $ minimum $s$ s.t. $\Pr[\sum_{j' \leq i} x_{j'} \langle \bfC_{1:i-1,j}, \bfC_{1:i-1, j'} \rangle > s] \leq \delta', x_{j'} \stackrel{i.i.d.}{\sim} Bern(p)$\newline \Comment{$s_{i,j}$ is a tail bound on the dot product of first $i-1$ entries of $\bfC \bfx$ and $\bfC_{1:i-1,j}$.}
\State $\epsilon_{i, j} = \frac{z \ltwo{\bfC_{1:i-1, j}}}{\sigma} + \frac{2s_{i,j}- \ltwo{\bfC_{1:i-1, j}}^2}{2\sigma^2}$ \newline \Comment{$\epsilon_{i,j}$ is a tail bound on the privacy loss of a participation in round $j$ after outputting first $i-1$ rounds}
\State $\tilde{p}_{i, j} = \frac{p \cdot \exp\left(\epsilon_{i, j}\right)}{p \cdot \exp\left(\epsilon_{i, j}\right) + (1-p)}$
\EndIf
\EndFor
\State \Return $\{\tilde{p}_{i,j}\}_{i \in [m], j \in [n]}$. 
\end{algorithmic}
\end{algorithm}
\caption{Generalization of \texttt{ProbabilityTailBounds}.}
\label{fig:cptb-subroutine-generalization}
\end{figure}

Note that if we want to analyze the privacy guarantee for an example in $D_i$, $i-1$, this is the same as analyzing the privacy guarantee for an example in $D_1$, if we use $\bfC$ with the first $i-1$ rows/columns cut off. Then, without loss of generality we only need to state a privacy analysis for examples in $D_1$ - to get a privacy guarantee that holds for all examples simultaneously, for each $D_i$ we can compute a privacy guarantee using the above reduction, and then take the worst of these. Further, for some classes of matrices, such as Toeplitz matrices, the examples in $D_1$ will have the worst privacy guarantee and thus it suffices to only analyze these examples.

We now show \texttt{Generalized-}\MMCC{}, given in \cref{fig:cptb-extension}, computes a valid privacy guarantee under $b$-min-sep sampling.

\begin{thm}\label{thm:mmcc-generalization}
Let $\epsilon$ be the output of \texttt{Generalized-}\MMCC{}. Then the matrix mechanism with matrix $\bfC$, $b$-min-sep sampling, sampling probability $p$, noise level $\sigma$ satisfies $(\epsilon, \delta_1 + \delta_2)$-DP (for examples in $D_1$).
\end{thm}
\begin{proof}
The algorithm is almost the same as \cref{thm:conditioning}, so we just need to justify the key differences. In particular, we need to justify (1) the deletion of columns, (2) the choice of $\tilde{p}^{(b)}_{i,j}$, and (3) the choice of $\bfC^{(b)}$.

(1) is justified by the proof of Theorem 4 in \cite{BandedMF}, which observes that the products of columns $j$ of $\bfC$ for which $j \pmod b \neq 1$ and the corresponding rows of $\bfx$ are independent of $D_1$, i.e. we can treat their products as public information. So it does not affect the privacy analysis to delete these rows/columns from $\bfC$/$\bfx$, and then view the resulting $\bfx$ as generated by i.i.d sampling every round with probability $bp$.

(2) and (3) are both justified if we use conditional composition over  sequential mechanisms corresponding to $b$ rows of $\bfC \bfx + \bfz$ instead of a single row. Each of these sequential mechanisms is a VMoG mechanism, which \cref{cor:scalardominance} allows us to reduce to the scalar PMoG mechanism defined in terms of $\bfC^{(b)}$ in \texttt{Generalized-}\MMCC{}. The probabilities $\tilde{p}^{(b)}$ are then valid to use in the conditional composition by the same argument as in \cref{thm:conditioning}, up to the adjustment to use $b \delta_1$ instead of $\delta_1$. This adjustment is valid, since we only use fraction $1/b$ of the values generated by \texttt{GeneralizedProbabilityTailBounds}, i.e. we are union bounding over $1/b$ as many ``bad'' events as in the original proof, so we can increase the allowed probability for each ``bad'' event by $b$ (which is implicitly done by increasing $\delta_1$ by $b$).
\end{proof}

One can verify that (i) for $b = 1$, \texttt{Generalized-}\MMCC{} is equivalent to \MMCC{}, and that (ii) if $\bfC$ is $b$-banded, \texttt{Generalized-}\MMCC{} is equivalent to the privacy analysis in \cite{BandedMF}.